% Part: incompleteness
% Chapter: incompleteness-provability
% Section: godels-paper

\documentclass[../../../include/open-logic-section]{subfiles}

\begin{document}

\olfileid{inc}{inp}{gop}

\olsection{গ্যোডেলের মূল গবেষণাপত্রের সঙ্গে তুলনা}

গ্যোডেলের ১৯৩১ সালের গবেষণাপত্রটি নিয়ে কিছু সময় কাটানো সার্থক। তার
ভূমিকায় আমরা এইমাত্র আলোচিত ধারণাগুলির রূপরেখা আছে। গবেষণাপত্রটির উপর
শুধু চোখ বোলালেও প্রতিটি পর্যায়ে কী ঘটছে সহজে বোঝা যায়: প্রথমে গ্যোডেল
বিধিবদ্ধ ব্যবস্থা $P$-কে বর্ণনা করেন (সংকেতবিন্যাস, স্বতঃসিদ্ধ,
প্রমাণবিধি); তারপর আদিম পুনরাবৃত্ত অপেক্ষক ও সম্বন্ধগুলি সংজ্ঞায়িত করেন;
তারপর দেখান যে $x B y$ আদিম পুনরাবৃত্ত এবং যুক্তি দেন যে আদিম পুনরাবৃত্ত
অপেক্ষক ও সম্বন্ধগুলি $\Th{P}$-তে উপস্থাপিত। এরপর তিনি উপরের মতোই
অসম্পূর্ণতা উপপাদ্য প্রমাণ করেন। ৩ নং বিভাগে দেখান, অপ্রমাণযোগ্য উক্তিটি
পাটিগণিতের ভাষার একটি বাক্য হিসেবে নেওয়া যায়। এখানেই $\beta$-লেমার
উৎপত্তি; পুনরাবৃত্ত অপেক্ষকগুলি $\Th{Q}$-তে উপস্থাপনযোগ্য দেখানোর সময়
অনুক্রম সামলাতে আমরাও এটি ব্যবহার করেছি। কোন ন্যূনতম স্বতঃসিদ্ধসমষ্টিই
যথেষ্ট, গ্যোডেল তা আলাদা করে দেখাননি; কিন্তু এখন আমরা জানি $\Th{Q}$-ই
কাজটি করে। সব শেষে ৪ নং বিভাগে তিনি দ্বিতীয় অসম্পূর্ণতা উপপাদ্যের একটি
প্রমাণের রূপরেখা দেন।

\end{document}
